\documentclass[12pt]{article}
\usepackage[utf8]{inputenc}
\usepackage[T1]{fontenc}
\usepackage[margin=1in]{geometry}
\usepackage{setspace}
\usepackage{xcolor}
\usepackage[colorlinks=true, urlcolor=blue!60!black, linkcolor=blue!60!black]{hyperref}
\setlength{\parindent}{2em}
\setlength{\parskip}{1.25em}
\renewcommand{\baselinestretch}{1.0}
\begin{document}
\begin{center}
{\LARGE \textit{Whole Woman's Health v. Jackson}, 595 U.S. 30 (2021)\par}
\vspace{0.8em}
{\large Case Brief --- Four-Page Edition\par}
\end{center}
\vspace{1em}

\section{Memory Jogger}

State sovereign immunity and traditional equitable limits under \textit{Ex parte Young} barred federal pre-enforcement suits against state judges, court clerks, and the Texas Attorney General to block Texas's private-enforcement abortion statute (S.B. 8), leaving only a narrow path against state executive licensing officials (\href{https://www.supremecourt.gov/opinions/21pdf/21-463_3ebh.pdf}{Slip Op.}).

\section{Detailed Case Facts}

Texas's S.B. 8, effective September 1, 2021, banned abortions after a fetal heartbeat was detected, about six weeks, in conflict with \textit{Roe} and \textit{Casey}. To thwart pre-enforcement review, it forbade state and local officials from enforcing the ban and authorized any private person to sue anyone who performed or aided an abortion, for at least \$10,000 per violation plus fees. It also tilted procedure: plaintiffs could choose among 254 counties, prevailing defendants could not recover fees, nonmutual preclusion was barred, and liability could apply retroactively. Abortion providers led by Whole Woman's Health sued under \S~1983, naming state judge Austin Reeve Jackson and clerk Penny Clarkston (as class representatives), Attorney General Ken Paxton, four state licensing officials, and private activist Mark Lee Dickson.

\section{Procedural History}

Judge Pitman in the Western District of Texas denied motions to dismiss. The Fifth Circuit refused to enjoin S.B. 8 and stayed proceedings. On September 1, 2021, the Supreme Court denied emergency relief 5--4. It granted certiorari before judgment, heard argument November 1, and decided the case December 10, 2021. On remand, the Texas Supreme Court held that Texas law barred the licensing officials from enforcing S.B. 8, ending the suit; \textit{Dobbs} then overruled \textit{Roe}.

\section{Judicial Votes}

Affirmed in part, reversed in part, and remanded.
\begin{itemize}
\item \textbf{Opinion of the Court:} Justice Gorsuch, joined by Justices Thomas, Alito, Kavanaugh, and Barrett (5--4 on clerks and Paxton; 9--0 on Judge Jackson and Dickson).
\item \textbf{Part II-C (plurality):} Gorsuch, Alito, Kavanaugh, and Barrett, with Roberts, Breyer, Sotomayor, and Kagan concurring, let claims against licensing officials proceed (8--1).
\item \textbf{Partial dissents:} Justice Thomas; Chief Justice Roberts, joined by Justices Breyer, Sotomayor, and Kagan; Justice Sotomayor, joined by Justices Breyer and Kagan.
\end{itemize}

\section{Holding}

\textit{Ex parte Young} is bounded by traditional equity. It permits pre-enforcement suits only against state executive officials with specific authority to enforce the challenged law or collateral disciplinary provisions. It does not allow injunctions against state judges adjudicating private suits, clerks docketing filings, or an attorney general lacking enforcement authority, and no court may enjoin unnamed private citizens or ``laws themselves.'' Licensing officials with apparent authority to discipline licensees survived the pleading stage; Dickson, who disclaimed any intent to sue, was dismissed for lack of standing. Gorsuch's assurances about alternative state-court avenues were dicta.

\section{Analysis of Opinions}

\textbf{Majority (Gorsuch).} \textit{Ex parte Young} is a narrow exception. Judges resolve disputes rather than enforce laws, and enjoining a state court or its machinery offends dual sovereignty; clerks performing ministerial docketing are not adverse to litigants, and no limiting principle would cabin clerk injunctions. Paxton had no statutory enforcement power over S.B. 8. In Part II-C, the plurality read S.B. 8's saving clause and Occupations Code \S~164.055 to leave the Medical Board power to discipline violators. Part III answered the dissent: there is no universal right to pre-enforcement review, federal defenses remain available in state court, and chill alone cannot override Article III.

\textbf{Thomas (in part).} Historical equity allows only a ``negative injunction'' against an officer with enforcement power who threatens imminent proceedings; S.B. 8 ``triply'' stripped all officials of authority. The Texas Supreme Court later agreed.

\textbf{Roberts, with Breyer, Sotomayor, and Kagan.} Paxton had civil-penalty authority, and because filing and docketing themselves inflicted the chill, clerks were proper defendants. Invoking \textit{Marbury} and \textit{United States v. Peters}, he warned that procedural artifice threatened the Court's role.

\textbf{Sotomayor, with Breyer and Kagan.} S.B. 8 deputized bounty hunters and turned state courts into a harassment weapon; likening it to Calhoun's nullification, she warned other States would copy it to chill gun, religious, or speech rights.

\section{Examples: Future Applications}

\begin{enumerate}
\item \textbf{Same-Side: Gun-Bounty Law.} A State bans possession of certain semiautomatic firearms through private \$15,000 suits and bars official enforcement. Suits against judges and clerks must be dismissed.
\item \textbf{Same-Side: Social Media Bounty.} A private-bounty law targets platforms' content moderation and strips the attorney general of enforcement power. He must be dismissed, and courts cannot enjoin unnamed private plaintiffs.
\item \textbf{Opposite-Side: Accreditation Revocation.} A separate code provision requires the education commissioner to revoke accreditation of violating schools. Under Part II-C, a suit against the commissioner may proceed.
\item \textbf{Opposite-Side: Threatening Private Enforcer.} An activist sends a theater owner demand letters threatening a bounty suit on Monday. Unlike Dickson, his imminent threat supports standing, and private persons have no immunity.
\item \textbf{Fence-Sitter: Gatekeeping Registrar.} A health registrar must stamp a ``Certificate of Compliance'' before any bounty suit. The registrar is an executive gatekeeper, but the role may be ministerial, like docketing.
\end{enumerate}

\section{Critique}

\textbf{Originalist and formalist views.} Defenders praise the Court's fidelity to traditional equity, Article III adversity, and the dignity of state judiciaries; federal courts lack roving power to enjoin court personnel or the public. Justice Thomas's reading of Texas law proved prescient.

\textbf{Progressive and remedies views.} Richard Fallon argued that rigid formalism truncated constitutional remedies and cheapened rights. Jon Michaels and David Noll warned of ``vigilante federalism,'' as States across the spectrum copied the model. Commentators also note tension with \textit{Pulliam v. Allen}. The case pits sovereign immunity and equitable restraint against the Supremacy Clause and effective remedies.

\section{Key Quotations}

\begin{quote}
``Consistent with historical practice, a federal court exercising its equitable authority may enjoin named defendants from taking specified unlawful actions. But under traditional equitable principles, no court may `lawfully enjoin the world at large,' \dots{} or purport to enjoin challenged `laws themselves.' ''
\end{quote}
\noindent --- Justice Gorsuch, Opinion of the Court.

\begin{quote}
``The clear purpose and actual effect of S. B. 8 has been to nullify this Court's rulings. . . . The nature of the federal right infringed does not matter; it is the role of the Supreme Court in our constitutional system that is at stake.''
\end{quote}
\noindent --- Chief Justice Roberts, concurring in the judgment in part and dissenting in part.

\begin{quote}
``As \textit{Ex parte Young} put it, `an injunction against a state court' or its `machinery' `would be a violation of the whole scheme of our Government.' ''
\end{quote}
\noindent --- Justice Gorsuch, Opinion of the Court.

\end{document}
